\section*{Chapitre \ref{ch:b2:living-soil} --- Le sol vivant}

\begin{solution}{exo:b2:living-soil:1}
O~: litière et matière organique partiellement décomposée, où commence la
décomposition. A~: \omterm{def:b2:living-soil:soil}{sol} minéral mêlé d'\omterm{def:b2:living-soil:soil}{humus}, sombre et à structure
grumeleuse, qui contient la plupart des racines, des organismes, des
nutriments et de l'eau. B~: l'horizon d'accumulation, où se déposent
l'argile, les oxydes de fer et les carbonates lessivés depuis le haut.
C~: matériau parental altéré, fragments de roche dans une matrice fine,
où se libère le nouveau \omterm{def:b2:living-soil:soil}{sol} minéral. R~: roche en place non altérée.
\end{solution}

\begin{solution}{exo:b2:living-soil:2}
Les charges négatives des surfaces de l'argile et de l'\omterm{def:b2:living-soil:soil}{humus} retiennent
des cations échangeables que les racines peuvent prélever en échange de
protons. Les grains de sable sont du quartz, non chargé et grossier,
offrant peu de surface par gramme et presque pas d'\omterm{def:b2:living-soil:soil}{humus}. Le chaulage
apporte du carbonate de calcium~: le calcium déplace les ions hydrogène
et aluminium qui saturent les sites d'échange d'un \omterm{def:b2:living-soil:soil}{sol} acide, le
carbonate neutralise l'acidité, et les sites sont de nouveau occupés par
un cation nutritif.
\end{solution}

\begin{solution}{exo:b2:living-soil:3}
Décomposeurs~: bactéries (\emph{Bacillus}) et champignons
(\emph{Trichoderma})~; brouteurs de microbes~: \omterm{def:b2:unicellular-diversity:protist}{protistes} (amibes) et
nématodes bactérivores~; détritivores~: \omterm{prop:b2:living-soil:community}{vers de terre}, diplopodes,
collemboles, cloportes~; prédateurs~: acariens, nématodes prédateurs,
chilopodes, carabes~; et taupes ou musaraignes au sommet. Les brouteurs,
dont le $\mathrm{C:N}$ est plus élevé que celui des bactéries dont ils se
nourrissent, excrètent l'azote excédentaire sous forme d'ammonium~: ils
minéralisent l'azote que les bactéries avaient immobilisé.
\end{solution}

\begin{solution}{exo:b2:living-soil:4}
Arbusculaires~: le champignon (un gloméromycète) pénètre dans les
cellules du cortex de la racine et y forme des arbuscules~; on les trouve
chez la plupart des herbacées et des plantes cultivées et chez les arbres
tropicaux. Ectomycorhiziennes~: le champignon (basidiomycètes et
ascomycètes, dont beaucoup de champignons à chapeau) engaine l'extrémité
de la racine et se développe entre les cellules sans y pénétrer~; on les
trouve chez les arbres tempérés et boréaux. Dans les deux cas, la plante
donne du sucre et le champignon donne du phosphate, de l'azote, de l'eau,
des oligo-éléments et une certaine protection contre les pathogènes.
\end{solution}

\begin{solution}{exo:b2:living-soil:5}
$L^{*} = I/k = 4/0.8 = \qty{5}{t/ha}$~; \omterm{def:b2:biogeochemical-cycles:reservoir}{temps de résidence} $1/k = 1.25$
an~; perte de \qty{90}{\%} au bout de $\ln 10/0.8 = 2.9$ ans.
\end{solution}

\begin{solution}{exo:b2:living-soil:6}
Pour \qty{100}{kg} de carbone consommé, les microbes en conservent
\qty{40}{kg} et ont besoin de $40/8 = \qty{5}{kg}$ d'azote. La paille en
fournit $100/80 =
\qty{1.25}{kg}$~: immobilisation de \qty{3.75}{kg} pour \qty{100}{kg}
de carbone. Le trèfle en fournit $100/15 = \qty{6.7}{kg}$~: minéralisation
de \qty{1.7}{kg}.
\end{solution}

\begin{solution}{exo:b2:living-soil:7}
\qty{20}{t/ha} de terre par an passent par les vers. Les \qty{20}{cm}
supérieurs pèsent $0.2\times 10^{4}\times 1.3 = \qty{2600}{t/ha}$~:
130 ans pour un passage. Les dix tonnes par acre de Darwin représentent
\qty{25}{t/ha}, le même ordre de grandeur --- son «~tous les quelques ans~»
concernait la terre végétale superficielle, quelques centimètres, et non
toute la couche arable.
\end{solution}

\begin{solution}{exo:b2:living-soil:8}
Bactéries~: $10^{9}\times 10^{-12} = \qty{1}{mg}$ par gramme. Hyphes~:
volume $\pi(2\times 10^{-4})^{2}\times 10^{4} = \qty{1.26e-3}{cm^3}$,
masse \qty{1.4}{mg} par gramme. Par hectare ($2.6\times 10^{9}$ g)~:
\qty{2.6}{t} de bactéries et \qty{3.6}{t} de champignons.
\end{solution}

\begin{solution}{exo:b2:living-soil:9}
$H^{*} = 1.5/0.02 = \qty{75}{t/ha}$ de carbone. Labouré~: $1.5/0.04
= \qty{37.5}{t/ha}$~; perte de \qty{37.5}{t/ha}, dont la moitié a
disparu au bout de $\ln 2/0.04 = 17$ ans.
\end{solution}

\begin{solution}{exo:b2:living-soil:10}
Phosphate~: $\sqrt{2\times 10^{-13}\times 8.64\times 10^{6}} =
\qty{1.3}{mm}$ en une saison --- une racine située à \qty{1}{cm} de la
suivante atteint un huitième du \omterm{def:b2:living-soil:soil}{sol} qui les sépare, tandis que des hyphes
espacées de \qty{0.1}{mm} l'atteignent en totalité. Nitrate~: \qty{4}{cm},
plus que l'espacement des racines~: il vient de lui-même jusqu'à la racine
(et par \omterm{def:b2:biogeochemical-cycles:reservoir}{flux} de masse avec l'eau), si bien que les mycorhizes apportent
peu pour le nitrate et beaucoup pour le phosphate.
\end{solution}

\begin{solution}{exo:b2:living-soil:11}
Masse $0.25\times 10^{4}\times 1.3 = \qty{3250}{t/ha}$. Perte nette de
\qty{19}{t/ha} par an~: durée de vie de 170 ans. Carbone perdu $20\times
0.03 = \qty{0.6}{t/ha}$ par an. En semis direct, la perte nette est de
\qty{1}{t/ha}~: 3250 ans.
\end{solution}

\begin{solution}{exo:b2:living-soil:12}
La litière se renouvelle en un an ou deux, l'\omterm{def:b2:living-soil:soil}{humus} en cinquante ans,
l'horizon A minéral en milliers d'années (sa formation étant de
\qty{0.1}{mm} par an), et le sel s'accumule en quelques décennies et
s'évacue en quelques décennies s'il y a un drainage. Un agriculteur voit
la litière et la culture réagir en une saison, et l'\omterm{def:b2:living-soil:soil}{humus} au cours d'une
carrière (une baisse d'un tiers prend trente ans~; une récupération tout
autant)~; l'érosion du \omterm{def:b2:living-soil:soil}{sol} minéral est invisible à l'échelle d'une vie et
irréversible à l'échelle de dix~; le sel est visible en une génération et,
sans drainage, permanent. Les variables lentes du \omterm{def:b2:living-soil:soil}{sol} sont celles
qu'aucune comptabilité annuelle n'enregistre.
\end{solution}

\begin{solution}{pb:b2:living-soil:1}
\textbf{1.} $0.25\times 10^{4}\times 1.3 = \qty{3250}{t/ha}$~;
carbone $\qty{81}{t/ha}$.
\textbf{2.} $81/12 = \qty{6.8}{t/ha}$ d'azote.
\textbf{3.} Bactéries \qty{1}{mg/g}, soit $\qty{3.25}{t/ha}$~; hyphes
$\pi(2\times 10^{-4})^{2}\times 2\times 10^{4}\times 1.1 =
\qty{2.8}{mg/g}$, soit \qty{9}{t/ha}~: biomasse microbienne d'environ
\qty{12}{t/ha}.
\textbf{4.} Environ \qty{14}{t/ha} de matière vivante (masse fraîche),
quelque \qty{17}{\%} du carbone de l'\omterm{def:b2:living-soil:soil}{humus} --- quelques pour cent de
celui-ci en carbone, puisque les organismes sont surtout constitués
d'eau.
\textbf{5.} Carbone microbien d'environ la moitié de \qty{12}{t}, soit
\qty{6}{t}~; en se renouvelant deux fois par an avec \qty{40}{\%}
retenus, les microbes consomment
$2\times 6/0.4 = \qty{30}{t}$ de carbone et en respirent \qty{18}{t/ha},
soit \qty{66}{t} de $\mathrm{CO_2}$ (une estimation grossière~: la
rétention et le renouvellement sont approximatifs).
\textbf{6.} Trois fois la production primaire nette de la culture, de
\qty{5}{t}~: l'estimation est trop élevée, car les hypothèses (toute la
biomasse se renouvelant deux fois, aucune fraction dormante) sont
généreuses~; la \omterm{prop:b2:living-soil:humus}{respiration du sol} mesurée sous une culture est de l'ordre de la
production de la culture elle-même, plusieurs tonnes de carbone par
hectare et par an. Le \omterm{def:b2:living-soil:soil}{sol} respire autant que les plantes qu'il porte.
\textbf{7.} Carbone $5\times 0.45 = \qty{2.25}{t}$~; azote $5\times
0.03 = \qty{150}{kg}$~; $\mathrm{C:N} = 15$.
\textbf{8.} Il minéralise~: les microbes ont besoin de $2250\times 0.4/8 =
\qty{112}{kg}$ d'azote et le résidu en contient 150, si bien que
\qty{38}{kg/ha} sont libérés en net.
\textbf{9.} $1 - e^{-0.5} = \qty{39}{\%}$ au bout de 3 mois~; $1 - e^{-2}
= \qty{86}{\%}$ au bout d'un an.
\textbf{10.} $0.39\times 38 = \qty{15}{kg/ha}$.
\textbf{11.} La culture a besoin de \qty{150}{kg}, pour l'essentiel dans
ses deux premiers mois~; le résidu en fournit \qty{15}{kg} en trois mois
et \qty{38}{kg} au total --- un complément, non un approvisionnement. Sa
valeur tient aussi à l'azote que le trèfle a fixé et que contiennent les
\qty{150}{kg} du résidu lui-même, dont l'essentiel entre dans l'\omterm{def:b2:living-soil:soil}{humus} et
revient au fil des années.
\textbf{12.} Paille~: carbone \qty{2250}{kg}, azote $2250/80 =
\qty{28}{kg}$~; les microbes ont besoin de \qty{112}{kg}, si bien que
\qty{84}{kg/ha} sont immobilisés aux dépens du \omterm{def:b2:living-soil:soil}{sol} --- plus de la moitié
des besoins de la culture. L'agriculteur doit apporter de l'azote avec la
paille, ou l'enfouir bien avant le semis, ou la laisser en surface, où
elle se décompose lentement.
\textbf{13.} Labour~: $1.2/0.025 = \qty{48}{t/ha}$~; semis direct~:
$1.2/0.015 = \qty{80}{t/ha}$.
\textbf{14.} $H(t) = 80 - 32\,e^{-0.015t}$~: gain de \qty{4.5}{t/ha}
au bout de 10 ans, de \qty{17}{t/ha} au bout de 50.
\textbf{15.} $\mathrm{d}H/\mathrm{d}t = 0.015\times 32 =
\qty{0.48}{t/ha}$ de carbone la première année, soit \qty{1.8}{t} de
$\mathrm{CO_2}$.
\textbf{16.} $0.48\times 1.5\times 10^{9} = \qty{0.7}{GtC}$ la
première année, \qty{7}{\%} des émissions fossiles~; la cinquantième
année, la vitesse est tombée à $0.48\,e^{-0.75} = \qty{0.23}{t/ha}$,
soit \qty{0.34}{GtC}.
\textbf{17.} Le gain est l'approche d'un nouvel \omterm{def:b2:biogeochemical-cycles:reservoir}{état stationnaire}~: il
tend vers zéro à mesure que $H$ atteint \qty{80}{t/ha}, et la totalité du
gain retourne à l'air en quelques décennies si le champ est de nouveau
labouré --- un puits dans le \omterm{def:b2:living-soil:soil}{sol} est un transfert unique et réversible,
non une compensation permanente d'émissions qui se poursuivent.
\textbf{18.} $k_h = 0.02$~: $H^{*} = \qty{60}{t/ha}$, soit une perte de
\qty{20}{t/ha}, plus que le gain des cinquante premières années du semis direct~: le réchauffement peut
défaire ce qu'a fait la gestion.
\textbf{19.} Labour~: $15 - 0.5 = \qty{14.5}{t/ha}$ par an~; semis
direct~: \qty{1}{t/ha}.
\textbf{20.} $3250/14.5 = 220$ ans~; $3250/1 = 3250$ ans.
\textbf{21.} Carbone $15\times 0.025 = \qty{0.38}{t/ha}$~; azote
$0.38/12 = \qty{31}{kg/ha}$ par an.
\textbf{22.} La décomposition de l'\omterm{def:b2:living-soil:soil}{humus} à l'\omterm{def:b2:biogeochemical-cycles:reservoir}{état stationnaire} du labour
vaut $k_hH^{*} =
I = \qty{1.2}{t/ha}$ par an~; l'érosion y ajoute encore un tiers, mais,
contrairement à la décomposition, elle n'est pas compensée par l'apport
--- elle vide le stock.
\textbf{23.} En partie. Le carbone érodé enfoui dans un sédiment ou un
\omterm{def:b2:biogeochemical-cycles:reservoir}{réservoir} peut être protégé de la décomposition pendant longtemps (un
puits par enfouissement)~; mais les agrégats sont brisés pendant le
transport, une grande partie du carbone est oxydée en chemin, et la
fertilité perdue du champ est remplacée par des engrais dont la
fabrication émet du carbone. Le signe global de l'effet de l'érosion sur
le carbone reste débattu.
\textbf{24.} Litière $1/k = 0.5$ an~; \omterm{def:b2:living-soil:soil}{humus} $1/k_h = 40$ ans sous le
labour et 67 en semis direct~; \omterm{def:b2:living-soil:soil}{sol} minéral $3250/0.5 =
6500$ ans au regard de la formation, et 220 ans au regard de la perte
sous le labour.
\textbf{25.} Stock de carbone de \qty{81}{t/ha} dans l'horizon A~; durée
de vie de 220 ans sous le labour, de 3250 ans en semis direct~; le résidu
de trèfle fournit environ \qty{38}{kg/ha} d'azote minéral, dont 15 dans
les trois premiers mois.
\end{solution}
